\section{The exact interface retained from the corpus}\label{sec:interface}

The normalized input is a bipartite Boolean tensor network
\[
 Z=\sum_{x\in\bits^E}\prod_{v}T_v(x|_{\delta(v)})
\]
with at most \(s\) edges, polynomially many vertices, maximum degree at
most \(\lceil s^{9/10}\rceil\), and description length \(\mathcal L\).
Its entries lie in
\[
  (\mathbb Z/2^{P_0}\mathbb Z)[V]/(V^\ell),\qquad
  1\le P_0,\ell\le s.
\]
Entry evaluation takes \((s+\mathcal L)^{O(1)}2^{O(s^{9/10})}\) time
and polynomial space. The extension ring, immutable creation-time
representatives, and exact interpolation conventions are unchanged.
In particular, canonical lifting is not assumed to preserve ring
operations \cite{source,asymmetric,pairwise}.

At an active call let \(M\) be the current wire count, \(r\) the fixed
group alphabet, \(u\) the retained valuation range, and \(L=2^a\) the
packet length. At a retained valuation \(i<u\) put \(D_i=2^{a-i}\).
The imported sharpened estimates give working precision
\[
 p_t\le(6+o(1))MD_i/u,
\]
and child wire count \((2/r+o(1))M\). The full-block Cauchy modes have
weights \(a+1-i+4j-2\pc(j)\), with at most
\((2/r^2+o(1))M/D_i\) families per directed pair at each index.
These are inherited lemmas, not new assertions proved by the finite
checker \cite{joint,pairwise}.

Fix an exact allocation of the weight budget among directed pairs.
Let \(\mathcal I_{gh}\) be the finite index list for a pair, skipping
allocations with an empty list. Its contribution is exactly
\begin{equation}\label{eq:exactmode}
 \sum_{\xi_{gh}\in\mathcal I_{gh}}
       \prod_{g=1}^r Z_g((\xi_e)_{e\ni g}).
\end{equation}
Each \(Z_g\) is a real recursive child contraction, not a unit-cost
oracle. For fixed \(r\), allocations and their bookkeeping contribute
polynomial factors. Exact knapsack counting and ranking give the
cardinalities and indices used below.

Write
\[
 h_{gh}=\log_2|\mathcal I_{gh}|,\qquad
 g_{\{g,h\}}=h_{gh}+h_{hg},\qquad
 K(M,u)=\kappa M/\sqrt u,\quad
 \kappa=\frac{\pi\sqrt2}{\ln2}.
\]
The common-tilt mode estimate implies, for every fixed nonnegative vector
\(v\) on unordered pairs,
\begin{equation}\label{eq:entropy}
 v\cdot g\le
 \left(\frac{\sqrt2\,\norm v}{r}+o(1)\right)K(M,u).
\end{equation}
The \(\sqrt2\) converts directed to unordered pair coordinates. The
error is uniform over the finite family of support vectors and sums
used in this paper.

The available asymmetric packing gates, with current side degrees
\(R,C\) and \(Q=r^a\), are
\begin{equation}\label{eq:gates}
 QR/M=o(1),\qquad
 Q^2C(\log s)^2/M=o(1),\qquad
 LQ^2/M=o(1).
\end{equation}
The retained side stays under the original degree bound; the fresh side
below the root has degree at most \(2^{a_{\rm previous}+1}\).

\section{A sharper rectangular operation bound}\label{sec:backend}

The cited matrix multiplication result proves
\(\alpha\ge0.321334\), \(\omega\le2.371552\), and the interior
rectangular bounds
\(\omega(1,0.55,1)\le2.066134\) and
\(\omega(1,0.60,1)\le2.092631\) \cite[Table 1]{matrix}.
Choose fixed rational values strictly inside the permitted bounds:
\[
 \alpha_0=\frac{321333}{10^6},\qquad
 \omega_0=\frac{2371553}{10^6},\qquad
 \eps=\frac1{10^6}.
\]
There are fixed algorithms of rectangular exponent \(2+\eps\) at aspect
\(\alpha_0\), exponents \(2.066134+\eps\) and \(2.092631+\eps\)
at aspects \(0.55\) and \(0.60\), and square exponent \(\omega_0\).
Every selected exponent has a fixed positive slack above its cited bound.
We do not claim these matrix multiplication results as new or as the
latest available matrix multiplication bounds.

Let \(f(k)\) be the line through
\((11/20,2066135/10^6)\) and \((3/5,2092632/10^6)\).
Tensor products of the two rectangular algorithms realize every rational
point of this segment. Put \(c=(\omega_0-1)/2=1371553/2000000\).
Intersect \(f(k)\) with \(1+c(1+k)\). Exact arithmetic gives
\[
 k_*=\frac{59261}{103891},\qquad
 b_*=\frac{26972100941}{12986375000}
     =f(k_*)=1+c(1+k_*).
\]
Here \(11/20<k_*<3/5\), and \(b_*\) exceeds the interpolated
\emph{published} table bound by exactly \(\eps\).
Thus the anchor has a genuine fixed-algorithm interpretation.

For sorted logarithmic dimensions \(x\ge y\ge z\ge0\), define
\begin{align}
 d&=1+\eps/2=\frac{2000001}{2000000},\notag\\
 w&=\frac{b_*-2-\eps}{k_*-\alpha_0}
     =\frac{7994703637}{25877393297},\notag\\
 v&=1+\eps-\alpha_0w
     =\frac{2913557133825647}{3234674162125000},\notag\\
 c&=\frac{1371553}{2000000}. \label{eq:backendcoeff}
\end{align}
In particular \(1>v>c>w>0\). The operation exponent is bounded by
\begin{equation}\label{eq:rectcost}
 F(x,y,z)=\max\{d(x+y),\ x+vy+wz,\ x+c(y+z)\}.
\end{equation}

\begin{lemma}[Rectangular bound]\label{lem:rect}
Formula~\eqref{eq:rectcost} pays for the matrix operation and its
input/output arrays, up to polynomial factors.
\end{lemma}
\begin{proof}
If \(y=0\), then \(z=0\) and the bound is immediate. Otherwise put
\(k=z/y\).
For \(\alpha_0\le k\le k_*\), tensoring the algorithms at
\(\alpha_0\) and \(k_*\) gives exponent
\(2+\eps+w(k-\alpha_0)\) in the balanced outer dimension \(y\).
Blocking the dimension \(x\) adds \(x-y\), giving \(x+vy+wz\).
For \(k_*\le k\le1\), the chord from \((k_*,b_*)\) to
\((1,\omega_0)\) is \(1+c(1+k)\); blocking gives \(x+c(y+z)\).
Both branches are bounded by the maximum in \eqref{eq:rectcost}.
Integer roundings of tensor powers affect only fixed factors.
If \(k<\alpha_0\), padding gives
\(x+y+\eps y\le d(x+y)\).
The largest input or output array has exponent at most \(x+y\), covered
by the first branch. The bilinear tensor symmetries permit sorting the
three logarithmic dimensions. Because \(1>v>w\) and \(1>c\), each
branch is the maximum of its coefficient permutations on an unsorted
triple. No exchange of minima and maxima is needed.
\end{proof}

We retain the corpus's coordinate/CRT bridge from fixed bilinear algorithms
over characteristic zero to the coefficient ring: lift coordinates to
bounded integers, use nonexceptional odd primes, reconstruct exact
products, and reduce modulo the required power of two. Polynomial rank,
precision, and logarithms of the dimensions make this a polynomial
overhead. Fixed bilinear recursions may process one child product at a time,
with polynomial space overhead relative to the largest input/output array.
The bridge is an imported implementation lemma \cite{rectangular};
the finite checker does not execute fast matrix multiplication.

\section{Twelve plans from an affine plane}\label{sec:plans}

Label nine group vertices by \((x,y)\in\F_3^2\), encoded as \(3x+y\),
and add vertex \(9\). Use the four linear forms
\[
 x,\quad y,\quad x+y,\quad x+2y.
\]
For each form \(a x+b y\), its three lines are
\(L_c=\{(x,y):ax+by=c\}\), \(c\in\F_3\).
For each of the three choices \(k\), make the ordered partition
\begin{equation}\label{eq:affineplans}
 (L_k\cup\{9\})\mid L_{c_1}\mid L_{c_2},
 \qquad c_1<c_2,\quad \{c_1,c_2\}=\{0,1,2\}\setminus\{k\}.
\end{equation}
Within each line the integer labels are increasing, and \(9\) is appended
to its block. This specifies twelve plans, each with block sizes \(4,3,3\).
Appendix~\ref{sec:plantable} gives the literal lists.

To build a block with ordered list \(B\), split at
\(\lfloor|B|/2\rfloor\), recursively build its two halves, and merge their
tables. For a merge with disjoint parts \(A,B\) and outside set \(O\),
the three index classes are the pairs \(AB,AO,BO\). The final operation
contracts the three completed block tables.
There are seven internal merges and one final triangle per plan.
The singleton tables are computed recursively and charged separately.

\subsection{Integer support vectors}

The sorted cost \eqref{eq:rectcost} is exactly the maximum of twelve linear
forms on the unsorted triple of dimension logarithms:
the three permutations of \((d,d,0)\), six of \((1,v,w)\), and three of \((1,c,c)\).
This follows from the rearrangement inequality.
The common integer scale is
\[
 \mathsf L=51754786594000000.
\]
The generating integer triples are
\begin{equation}\label{eq:integertriples}
 \mathsf L(d,d,0),\qquad
 \mathsf L(1,v,w),\qquad
 \mathsf L(1,c,c),
\end{equation}
with exact rational coefficients \eqref{eq:backendcoeff}. Each is integral;
the checker constructs the scale by least common multiples.
Assign the coefficient for a class to every unordered group edge in that
class, and zero to edges within a part. Each plan \(j\) has an indexed menu
\(\V_j\) of \(8\cdot12=96\) support vectors in
\(\mathbb Z_{\ge0}^{45}\); repeated vectors, if any, retain their indices.
Define its logarithmic local cost
\[
 F_j(g)=\mathsf L^{-1}\max_{v\in\V_j}v\cdot g.
\]
This is the maximum cost of all its matrix operations after the singleton
caches are available.

\subsection{One consistent choice per plan}

\begin{lemma}[Compatible-support bound]\label{lem:norm}
For \(q\) fixed plan menus, suppose
\[
 \left\|\sum_{j=1}^q v_j\right\|_2^2\le B
 \quad\hbox{for every }v_j\in\V_j.
\]
Selecting the plan of minimum \(F_j(g)\) has local time and space
exponents at most \((\beta+o(1))K(M,u)\), where
\[
 \beta^2=\frac{2B}{\mathsf L^2q^2r^2}.
\]
\end{lemma}
\begin{proof}
For the given \(g\), choose a maximizing support \(v_j\) from each menu.
Then
\[
 \min_jF_j(g)\le\frac1q\sum_jF_j(g)
  =\frac1{\mathsf L q}\left(\sum_jv_j\right)\cdot g
  \le\left(\frac{\sqrt{2B}}{\mathsf L qr}+o(1)\right)K
\]
by \eqref{eq:entropy}. Every choice belongs to one fixed finite set, so
the error bound is uniform. The input/output term in
\eqref{eq:rectcost} pays for every local table; summing a fixed number
of table sizes changes only a polynomial factor.
\end{proof}

The preceding pair certificate upper-bounded this norm by independent
maxima for each pair of menus. We instead certify its exact global maximum:
\begin{theorem}[Exact finite affine certificate]\label{thm:finite}
For the twelve menus specified above,
\begin{equation}\label{eq:exactmax}
 \max_{v_j\in\V_j}
       \left\|\sum_{j=1}^{12}v_j\right\|_2^2
 =6002555704965940465476145427912572416.
\end{equation}
Consequently the value \(\beta^2=B_*\) in
\eqref{eq:mainconstant} is valid.
\end{theorem}

\subsection{Proof certificate and its completeness}

We give the mathematical verification rule explicitly. The supplied
programs instantiate it using integers; their searches are finite proof
computations, not floating-point estimates of the maximum.

For choices \(v_i,v_j\) define the pair potential
\[
 P_{ij}(v_i,v_j)
 =\norm{v_i}^2+\norm{v_j}^2+2(q-1)v_i\cdot v_j.
\]
Expansion gives the exact identity
\begin{equation}\label{eq:energyidentity}
 \sum_{i<j}P_{ij}(v_i,v_j)
 =(q-1)\left\|\sum_i v_i\right\|_2^2.
\end{equation}
Every diagonal term occurs \(q-1\) times. For \(q=12\), the comparison
threshold is eleven times the proposed maximum in \eqref{eq:exactmax}.

Start with all 96 choices in each domain. Given current domains \(D_j\),
an upper bound for the pair energy conditional on \(v_i=a\) is
\begin{equation}\label{eq:delete}
 U_i(a)=
 \sum_{\substack{j<k\\i\notin\{j,k\}}}
       \max_{b\in D_j,c\in D_k}P_{jk}(b,c)
 +\sum_{j\ne i}\max_{b\in D_j}P_{ij}(a,b),
\end{equation}
where the last expression is interpreted symmetrically. If this is below
the threshold, no exceeding tuple can use \(a\). Delete all such choices
and repeat. Simultaneous deletions are sound: any exceeding tuple uses
only choices surviving the previous round, and its conditional upper
bound exceeds the threshold at every coordinate.

After elimination, fix a prefix \(v_1,\ldots,v_d\). Let \(E_d\) be its
exact internal pair energy, and for each unfixed \(j\) put
\(u_j(b)=\sum_{i\le d}P_{ij}(v_i,b)\). Then
\begin{equation}\label{eq:branchupper}
 U(\text{prefix})=
 E_d+\sum_{j>d}\max_{b\in D_j}u_j(b)
       +\sum_{d<j<k}\max_{b\in D_j,c\in D_k}P_{jk}(b,c)
\end{equation}
is an upper bound on every completion. If it is at most the threshold,
the entire branch is certified. Otherwise branch on every remaining
choice of the next coordinate. The depth is at most twelve.
An unpruned complete tuple is checked exactly. Induction on the number
of remaining coordinates proves that successful completion certifies
every tuple.

Here the scale and triples are \eqref{eq:integertriples}, pair coordinates
are in lexicographic order, operation traversal is postorder, and the twelve
triples are sorted lexicographically. These conventions uniquely determine
support indices. The witness takes index \(88\) in every plan.
Its directly computed norm equals \eqref{eq:exactmax}, establishing the
lower inequality. The upper certificate has:

\begin{center}
\begin{tabular}{@{}lr@{}}
\toprule
Certificate quantity & Exact count or value\\
\midrule
Support vectors & \(1152\)\\
Pair choices computed & \(66\cdot96^2=608256\)\\
Surviving choices & Two domains of \(25\); ten domains of \(21\)\\
Remaining Cartesian choices covered & \(25^2\,21^{10}\)\\
Branch-and-bound nodes visited & \(303513\)\\
Maximum norm witness & All twelve support indices equal \(88\)\\
\bottomrule
\end{tabular}
\end{center}

All remaining branches are certified before depth twelve; the equality
witness is evaluated separately. Zero complete leaves visited therefore
does not mean the Cartesian product was ignored.
The maximum discarded conditional upper bound is
\[
 66006990952377090594705594189055735809,
\]
strictly below the threshold
\[
 66028112754625345120237599707038296576.
\]
The first checker constructs incremental conditional sums. The independent
checker reconstructs supports differently and recomputes each prefix
bound directly. Both use arbitrary-precision integers and explicit
exception checks. Substitution into Lemma~\ref{lem:norm} gives \(B_*\).
The equality witness has a compact normalized expression:
\[
 B_*=\frac{47+116c+128c^2}{600}
     =\frac{5835847444309}{18750000000000}.
\]
The expression counts the edges in the sum of the twelve final-operation
supports \((c,1,c)\). Its equality with the integer witness is checked
exactly; the global upper bound still requires the compatibility certificate.

\subsection{Exact plan selection}

Let \(N_{\{g,h\}}=|\mathcal I_{gh}||\mathcal I_{hg}|\). For plan \(j\)
compute the integer
\[
 Q_j=\max_{v\in\V_j}\prod_e N_e^{v_e}.
\]
Minimizing \(Q_j\) is exactly equivalent to minimizing \(F_j(g)\).
The number of factors and their exponents are fixed constants, and the
cardinalities have polynomial bit length. This is an exact polynomial-time
selector; it requires no oracle for comparing logarithms. The large
denominator makes this a theoretical implementation specification rather
than a claim of practical efficiency.

